\chapter{Theory of the Travelling-Wave Parametric Amplifier}
\label{chap:theory}

This chapter develops the physics the solver is built to compute. It proceeds
from the single Josephson element to the distributed line, derives the
coupled-mode equations and the standard gain expression, and then --- the point
of the chapter --- identifies exactly which assumptions in that derivation fail
at the operating points of interest, and therefore what a numerical method has
to do differently.

\section{The Josephson element as a nonlinear inductor}

\subsection{The Josephson relations}

A tunnel junction between two superconductors supports a dissipationless
current set by the gauge-invariant phase difference $\varphi$ across the
barrier,
\begin{equation}
  I = I_c \sin\varphi,
  \label{eq:dc-josephson}
\end{equation}
while the phase evolves under an applied voltage as
\begin{equation}
  \frac{\mathrm{d}\varphi}{\mathrm{d}t} = \frac{2e}{\hbar}\,V
  = \frac{V}{\rphi},
  \qquad
  \rphi \equiv \frac{\hbar}{2e} = \frac{\fluxquantum}{2\pi}.
  \label{eq:ac-josephson}
\end{equation}

It is convenient to work not with $\varphi$ but with the \emph{branch flux}
\begin{equation}
  \psi(t) = \int_{-\infty}^{t} V(t')\,\mathrm{d}t' = \rphi\,\varphi(t),
\end{equation}
because $\psi$ is the natural variable of the node-flux circuit formalism used
throughout this work (\cref{chap:circuit}). In this variable the branch law
\eqref{eq:dc-josephson} reads
\begin{equation}
  \boxed{\;
  i_J(\psi) = I_c \sin\!\left(\frac{\psi}{\rphi}\right).
  \;}
  \label{eq:branch-law}
\end{equation}

\subsection{Nonlinear inductance and the expansion parameter}

Differentiating \eqref{eq:branch-law} gives the differential inverse inductance
\begin{equation}
  \frac{1}{L(\psi)}
  \;=\;
  \frac{\partial i_J}{\partial \psi}
  \;=\;
  \frac{I_c}{\rphi}\cos\!\left(\frac{\psi}{\rphi}\right)
  \;\equiv\;
  \gamma(\psi).
  \label{eq:gamma-def}
\end{equation}
The quantity $\gamma$ recurs constantly: it is the junction \emph{tangent}, and
its time dependence under a pump is the entire physical origin of parametric
gain. At zero flux it takes its maximum value
\begin{equation}
  \frac{1}{L_J} = \frac{I_c}{\rphi},
  \qquad\text{i.e.}\qquad
  L_J = \frac{\rphi}{I_c},
\end{equation}
which is the small-signal Josephson inductance. Expanding
\eqref{eq:branch-law},
\begin{equation}
  i_J(\psi) = \frac{\psi}{L_J}
  \left[
    1 - \frac{1}{6}\left(\frac{\psi}{\rphi}\right)^{2}
      + \frac{1}{120}\left(\frac{\psi}{\rphi}\right)^{4}
      - \cdots
  \right].
  \label{eq:branch-expansion}
\end{equation}

Two facts about \eqref{eq:branch-expansion} determine the character of the whole
problem:

\begin{enumerate}
\item \textbf{The leading nonlinearity is cubic, not quadratic.} An unbiased
      junction has an odd branch law, so $i_J(-\psi)=-i_J(\psi)$ and all even
      terms vanish identically. A cubic nonlinearity mediates
      \emph{four}-wave mixing (4WM): three photons in, one out. Obtaining
      three-wave mixing (3WM) requires deliberately breaking the symmetry, by
      a DC current bias, by a flux-biased asymmetric loop (SNAIL, SQUID), or
      by an asymmetric junction array.

\item \textbf{The expansion parameter is $\psi/\rphi$, and it is not small at
      the operating point.} A device delivering \SI{20}{\dB} of gain over a
      few thousand cells typically runs at $\psi/\rphi \sim 0.3$--$0.6$ per
      junction. At $\psi/\rphi = 0.5$ the quintic term in
      \eqref{eq:branch-expansion} is already a percent-level correction to the
      cubic one, and the higher terms are not negligible in the pump waveform
      even where they are negligible in the signal response. This is the first
      concrete reason why the solver evaluates \eqref{eq:branch-law} exactly in
      the time domain rather than truncating \eqref{eq:branch-expansion}.
\end{enumerate}

\subsection{Flux-tunable and composite elements}
\label{sec:composite-elements}

Several element types generalise \eqref{eq:branch-law}. All of them enter the
solver through the same interface --- a branch current $i(\psi)$ and its
tangent $\partial i/\partial\psi$ --- so the numerics of
\cref{chap:pumphb} are unchanged by the choice.

\paragraph{DC SQUID.} Two junctions in a loop threaded by $\Phi_{\mathrm{ext}}$
behave, for a symmetric pair, as a single junction with flux-tunable critical
current
\begin{equation}
  I_c(\Phi_{\mathrm{ext}}) = 2I_{c0}
  \left|\cos\!\left(\pi\frac{\Phi_{\mathrm{ext}}}{\fluxquantum}\right)\right|.
\end{equation}
The branch law is still odd, so a symmetric SQUID gives 4WM; asymmetry between
the two junctions introduces an even component and hence 3WM.

\paragraph{SNAIL.} The Superconducting Nonlinear Asymmetric Inductive eLement
places $N$ large junctions in parallel with one junction of scaled critical
current $r I_{c0}$, in a flux-biased loop. Its branch law is
\begin{equation}
  i(\psi) = I_{c0}\left[
    r\,\sin\!\left(\frac{\psi}{\rphi}\right)
    + \sin\!\left(\frac{\psi - \psi_{\mathrm{ext}}}{N\rphi}\right)
  \right],
  \label{eq:snail-law}
\end{equation}
for the convention in which the external flux threads the large-junction arm;
the alternative convention exchanges which term carries $\psi_{\mathrm{ext}}$.
Away from $\psi_{\mathrm{ext}}=0$ this is neither odd nor even, so it supports
3WM with a tunable ratio of quadratic to cubic nonlinearity --- the design
freedom that makes the SNAIL attractive.

Equation~\eqref{eq:snail-law} does not vanish at $\psi=0$. The element sits at a
\emph{static equilibrium} $\psi_{\mathrm{eq}}$ solving $i(\psi_{\mathrm{eq}})=0$,
and the physically meaningful expansion is about that point. Writing
$\psi = \psi_{\mathrm{eq}} + \delta\psi$ and expanding,
\begin{equation}
  i(\psi_{\mathrm{eq}} + \delta\psi)
  = I_{c0}\left[
      g_1 \frac{\delta\psi}{\rphi}
    + g_2 \left(\frac{\delta\psi}{\rphi}\right)^{2}
    + g_3 \left(\frac{\delta\psi}{\rphi}\right)^{3}
    + \cdots \right],
\end{equation}
with
\begin{align}
  g_1 &= r\cos\varphi_s + \tfrac{1}{N}\cos\varphi_l, \\
  g_2 &= -\tfrac{1}{2}\!\left(r\sin\varphi_s + \tfrac{1}{N^{2}}\sin\varphi_l\right), \\
  g_3 &= -\tfrac{1}{6}\!\left(r\cos\varphi_s + \tfrac{1}{N^{3}}\cos\varphi_l\right),
\end{align}
where $\varphi_s = \psi_{\mathrm{eq}}/\rphi$ and
$\varphi_l = (\psi_{\mathrm{eq}}-\psi_{\mathrm{ext}})/(N\rphi)$. Only the root
with $g_1>0$ is a stable equilibrium (positive effective inductance); the
solver selects it by that criterion.

\begin{implnote}[shifted SNAIL branch law]
\code{EffectiveSnailBranchLaw} in \file{core/nonlinear.py} implements
\eqref{eq:snail-law} \emph{shifted to its solved static equilibrium}: the flux
argument passed in is the deviation $\delta\psi$, and the class adds
$\psi_{\mathrm{eq}}$ internally through \code{\_absolute\_flux}. The equilibrium
itself is found by \code{snail\_taylor\_coefficients}, which brackets every root
of $r\sin x + \sin((x-\psi_{\mathrm{ext}}/\rphi)/3)$ on a
$6\pi$-wide grid, refines with \code{brentq}, and returns the first root with
$g_1>0$.

Because the element is presented already shifted, its small-signal slope
$g_1 I_{c0}/\rphi$ is carried \emph{entirely} by the nonlinear branch. Stamping
that slope into the linear stiffness matrix $\Kmat$ as well double-counts it.
That was a real defect in an earlier version of the SNAIL builder and it
destroyed the model's gain; see \cref{sec:k-vs-branch}.
\end{implnote}

\paragraph{Kinetic inductance.} A high-normal-state-resistance superconducting
film (NbTiN, NbN) has a current-dependent kinetic inductance
\begin{equation}
  L(I) = L_0\left[1 + \left(\frac{I}{I_*}\right)^{2} + \cdots\right],
\end{equation}
which is again even in $I$ and therefore gives 4WM. Kinetic-inductance TWPAs
(KITWPAs) avoid junction fabrication entirely at the cost of much higher pump
power. The solver's branch-law interface accommodates them, though the designs
studied here are junction-based.

\section{From lumped element to travelling wave}

\subsection{The nonlinear transmission line}

The canonical Josephson TWPA (JTWPA) is a ladder: a series chain of junctions
with a shunt capacitance $C_g$ to ground at each node, and a junction shunt
capacitance $C_J$ across each junction. Writing $\phi_j$ for the flux at node
$j$ and $a$ for the cell length, current conservation at node $j$ reads
\begin{equation}
  C_g\ddot{\phi}_j
  + C_J\!\left(2\ddot{\phi}_j - \ddot{\phi}_{j+1} - \ddot{\phi}_{j-1}\right)
  + I_c\!\left[
      \sin\!\frac{\phi_j-\phi_{j+1}}{\rphi}
    + \sin\!\frac{\phi_j-\phi_{j-1}}{\rphi}
    \right] = 0 .
  \label{eq:jtwpa-node}
\end{equation}

In the continuum limit $\phi_j \to \phi(x)$, retaining the leading nonlinearity
from \eqref{eq:branch-expansion}, one obtains
\begin{equation}
  \frac{\partial^{2}\phi}{\partial x^{2}}
  - \frac{1}{v^{2}}\frac{\partial^{2}\phi}{\partial t^{2}}
  + \frac{a^{2}C_J}{C_g}\frac{\partial^{4}\phi}{\partial x^{2}\partial t^{2}}
  \;=\;
  \frac{a^{2}}{6\rphi^{2}}
  \frac{\partial}{\partial x}\!\left[
    \left(\frac{\partial\phi}{\partial x}\right)^{3}
  \right],
  \label{eq:continuum}
\end{equation}
with phase velocity $v = a/\sqrt{L_J C_g}$ and characteristic impedance
$Z_c = \sqrt{L_J/C_g}$.

The linear part of \eqref{eq:continuum} gives the dispersion relation
\begin{equation}
  k(\omega) = \frac{\omega}{v}\,
  \frac{1}{\sqrt{1 - \omega^{2}/\omega_{J}^{2}}},
  \qquad
  \omega_J = \frac{1}{\sqrt{L_J C_J}},
  \label{eq:dispersion}
\end{equation}
where $\omega_J$ is the junction plasma frequency. The line is a low-pass
structure with cutoff at $\omega_J$; operating well below it keeps dispersion
weak, which as we will see is both necessary and problematic.

\subsection{Why the plasma frequency is held fixed}

Equation~\eqref{eq:dispersion} shows that $\omega_J$ sets the entire dispersion
profile. If one varies $L_J$ along the line --- to taper impedance, to apply a
deliberate profile, or through fabrication scatter --- and holds $C_J$ fixed,
then $\omega_J$ varies too and the dispersion becomes position dependent in a
way that is usually unintended.

\begin{implnote}[$C_J$ tracks $L_J$]
The circuit builders therefore derive the nominal junction capacitance from the
junction inductance so that $\omega_J$ is constant along the line:
\[
  C_J^{(j)} = C_J^{\mathrm{nom}}\,\frac{L_J^{\mathrm{nom}}}{L_J^{(j)}} .
\]
There is deliberately \emph{no} independent $C_J$ \emph{profile}. Independent
$C_J$ \emph{scatter} is supported and intentionally breaks the relation, because
fabrication scatter in the junction area affects $I_c$ and the capacitance
through different geometric dependences; that must not be ``corrected'' to track
$L_J$.
\end{implnote}

\section{Parametric mixing}

\subsection{Wave-mixing processes}

Write the field on the line as a sum of pump, signal and idler,
\begin{equation}
  \phi(x,t) = \tfrac{1}{2}\sum_{\mu\in\{p,s,i\}}
  A_\mu(x)\,e^{i(k_\mu x - \omega_\mu t)} + \mathrm{c.c.}
\end{equation}
Substituting into the nonlinear term of \eqref{eq:continuum} and collecting
resonant contributions gives the mixing processes the cubic nonlinearity
supports. The one that produces gain is
\begin{equation}
  \boxed{\;2\omega_p = \omega_s + \omega_i\;}
  \qquad\text{(degenerate 4WM)},
  \label{eq:4wm-energy}
\end{equation}
in which two pump photons are annihilated and a signal--idler pair is created.
If the symmetry is broken so that a quadratic term exists, the corresponding
3WM process is
\begin{equation}
  \omega_p = \omega_s + \omega_i
  \qquad\text{(3WM)}.
  \label{eq:3wm-energy}
\end{equation}

Both are phase-insensitive when $\omega_s\neq\omega_i$: the amplifier adds at
least half a photon of noise, and the idler carries away an equal photon flux.
The degenerate case $\omega_s=\omega_i$ is phase sensitive and can in principle
be noiseless in one quadrature, at the cost of amplifying only that quadrature.

\subsection{Coupled-mode equations}

Under the slowly-varying-envelope approximation --- $|\partial_x A| \ll |kA|$
--- and assuming an undepleted pump, the 4WM process gives
\begin{align}
  \frac{\mathrm{d}A_s}{\mathrm{d}x}
    &= i\kappa_s |A_p|^{2} A_s
     + i\kappa_c A_p^{2} A_i^{*} e^{i\Delta k\,x}, \label{eq:cme-s}\\[2pt]
  \frac{\mathrm{d}A_i^{*}}{\mathrm{d}x}
    &= -i\kappa_i |A_p|^{2} A_i^{*}
     - i\kappa_c^{*} (A_p^{*})^{2} A_s e^{-i\Delta k\,x}, \label{eq:cme-i}
\end{align}
with linear phase mismatch
\begin{equation}
  \Delta k = 2k_p - k_s - k_i .
  \label{eq:linear-mismatch}
\end{equation}
The $\kappa_{s,i}|A_p|^{2}$ terms are \emph{self- and cross-phase modulation}
(SPM/XPM): the strong pump shifts the refractive index seen by signal and idler.
The $\kappa_c$ term is the conversion that produces gain.

\subsection{Phase matching, and why 4WM does not phase match by itself}

Solving \eqref{eq:cme-s}--\eqref{eq:cme-i} with $A_s(0)=A_{s0}$, $A_i(0)=0$
gives exponential growth
\begin{equation}
  G(x) = \left|\frac{A_s(x)}{A_{s0}}\right|^{2}
       = 1 + \left(\frac{\kappa_c|A_p|^{2}}{g}\right)^{2}\sinh^{2}(gx),
  \label{eq:gain-cme}
\end{equation}
with the parametric gain coefficient
\begin{equation}
  g = \sqrt{\left(\kappa_c|A_p|^{2}\right)^{2}
            - \left(\frac{\Delta k_{\mathrm{tot}}}{2}\right)^{2}},
  \qquad
  \Delta k_{\mathrm{tot}} = \Delta k + (2\kappa_p - \kappa_s - \kappa_i)|A_p|^{2}.
  \label{eq:g-coefficient}
\end{equation}

Everything hinges on $\Delta k_{\mathrm{tot}}$. If it is small, $g$ is real and
the gain grows exponentially, $G\sim \tfrac14 e^{2gL}$. If
$|\Delta k_{\mathrm{tot}}| > 2\kappa_c|A_p|^{2}$, then $g$ becomes imaginary,
$\sinh\to\sin$, and the ``gain'' merely oscillates with position --- the
familiar $\mathrm{sinc}$-like ripple, with no net amplification.

The problem specific to 4WM is that the nonlinear contribution to
$\Delta k_{\mathrm{tot}}$ does not vanish. For a self-phase-modulation-dominated
line one finds $2\kappa_p - \kappa_s - \kappa_i \approx -\kappa_p$, so even a
perfectly dispersionless line ($\Delta k = 0$) has
\begin{equation}
  \Delta k_{\mathrm{tot}} \approx -\kappa_p |A_p|^{2} \neq 0 .
\end{equation}
The pump self-phase modulation \emph{itself} spoils the phase matching, and it
does so in proportion to the pump power --- exactly the knob one wants to turn
up. Left alone, a 4WM JTWPA saturates in gain at a few dB regardless of length.

3WM does not suffer this: the process is $\omega_p=\omega_s+\omega_i$ and the
leading nonlinear phase shift is quadratic in field rather than cubic, so the
pump-induced mismatch is far weaker. This is the main practical argument for
flux-biased 3WM devices.

\subsection{Dispersion engineering}

The standard fix for 4WM is to add a small, deliberately placed dispersive
feature that supplies a compensating linear mismatch. Two families are used:

\paragraph{Resonant phase matching (RPM).} Periodically loading the line with
LC resonators of frequency $\omega_r$ creates a narrow stopband. Just below
$\omega_r$ the phase velocity is strongly modified, so placing the pump there
gives a large tunable $\Delta k$ that can be set to cancel
$-\kappa_p|A_p|^{2}$. The cost is a narrow gap in the gain profile near
$\omega_r$ where the pump lives.

\paragraph{Periodic loading / photonic crystal.} Modulating the cell impedance
with period $d$ opens a bandgap at $k = \pi/d$. Operating the pump near the
band edge supplies the required dispersion without a discrete resonator.

Both introduce structure on a length scale comparable to the wavelength, which
means the continuum limit \eqref{eq:continuum} is no longer a good description
of the very feature the design relies on. This is a second concrete reason to
solve the discrete circuit \eqref{eq:jtwpa-node} directly.

\begin{caveat}[the dead band]
Dispersion-engineered 4WM devices have a frequency interval near the pump where
no useful gain exists. In the solver's frequency-resolved compression campaigns
this appears as the status \code{NO\_GAIN\_AT\_OPERATING\_POINT}. It is retained
as a distinct status rather than reported as an anomalously large or small
$P_{1\mathrm{dB}}$, because a compression point is undefined where there is no
gain to compress.
\end{caveat}

\section{Saturation}

\subsection{Pump depletion}

Equation~\eqref{eq:gain-cme} assumes an undepleted pump. That assumption fails
by construction at large signal power: the amplified signal and the generated
idler take their energy from the pump. Photon-number conservation
(Manley--Rowe, \cref{sec:manley-rowe}) requires that for every signal photon
added, one idler photon is created and --- in 4WM --- two pump photons are
destroyed.

A crude but instructive model retains only this bookkeeping. If $G_0$ is the
small-signal power gain and $P_s$, $P_p$ the signal and pump powers, then
\begin{equation}
  G(P_s) \simeq \frac{G_0}{1 + 2G_0 P_s/P_p}.
  \label{eq:depletion-model}
\end{equation}

\begin{implnote}[depletion-only model]
\code{depletion\_only\_model} in \file{multitone/compression\_curve.py}
implements \eqref{eq:depletion-model} exactly, returning \emph{linear} gain;
\code{depletion\_only\_gain\_db} converts to \si{\dB}. The CSV field
\code{compression\_model\_depletion\_only} carries the \si{\dB} value.
\end{implnote}

Equation~\eqref{eq:depletion-model} predicts $1$\,dB of compression when
$2G_0P_s/P_p = 10^{1/10}-1 \approx 0.2589$, i.e.
\begin{equation}
  P_{1\mathrm{dB}} \approx \frac{0.1294\,P_p}{G_0},
\end{equation}
so on a log--log plot $P_{1\mathrm{dB}}$ falls with a slope of exactly
$-1$\,dB per dB of gain, and the output-referred compression point
$P_{1\mathrm{dB}} + G_0 - 1$ is a gain-independent constant
($\approx P_p - \SI{9.87}{\dB}$, i.e.\ $\SI{1.302}{\dB}$ below
$P_p/8$ in the convention used here).

\begin{caveat}[the $-1$\,dB/dB slope is a diagnostic, not a bound]
Both the $-1$\,dB/dB slope and the constant output compression point are
predictions of one specific model, the pure-depletion model. They are useful as
a \emph{reference line}: a measured slope shallower than $-1$ indicates a
gain-independent limiting mechanism (something saturating that does not scale
with $G_0$), while a steeper slope indicates a super-linear one. Neither is a
physical bound and neither may be used as a pass/fail gate.
\end{caveat}

\subsection{Other saturation mechanisms}

Pump depletion is not the only route to compression, and in real devices it is
often not the dominant one:

\begin{itemize}
\item \textbf{Junction phase excursion.} As $\psi/\rphi$ grows, the tangent
      $\gamma$ of \eqref{eq:gamma-def} decreases and eventually changes sign. A
      junction driven past $\psi/\rphi = \pi/2$ has negative differential
      inductance over part of the cycle. The solver reports
      \code{min\_cosine} per junction (\code{junction\_diagnostics}) precisely
      to detect this.

\item \textbf{Higher-order mixing.} Once the signal is comparable to the pump,
      cascaded processes generate tones at $\omega_s \pm n\wp$ for large $n$,
      each draining power. This is what the multitone basis of
      \cref{chap:saturation} is designed to represent.

\item \textbf{Phase-mismatch drift.} Depletion changes $|A_p|^{2}$, which
      through \eqref{eq:g-coefficient} changes $\Delta k_{\mathrm{tot}}$ and
      hence detunes the device from its phase-matched point. This is a
      \emph{second-order} effect of depletion and can dominate the first-order
      power loss.

\item \textbf{Loss.} Line loss reduces the pump along the chain, so the
      effective $|A_p|^{2}$ in \eqref{eq:g-coefficient} is position dependent
      even without depletion.
\end{itemize}

None of these is captured by \eqref{eq:depletion-model}. Distinguishing them
requires the full nonlinear multitone solve of \cref{chap:saturation} together
with the spatial attribution diagnostics of \cref{sec:spatial}.

\section{Where the analytic theory stops}
\label{sec:where-theory-stops}

The coupled-mode picture developed above is the right way to \emph{think} about
a TWPA. It is not adequate to \emph{compute} one at the operating points of
interest. The specific failures, each of which motivates a corresponding feature
of the solver, are:

\begin{enumerate}
\item \textbf{Truncated nonlinearity.} \Cref{eq:cme-s}--\eqref{eq:cme-i} retain
      only the cubic term. At $\psi/\rphi\sim 0.5$ the pump waveform has
      several percent of third-harmonic content, and the higher odd harmonics
      are not negligible.
      \emph{Solver response:} evaluate \eqref{eq:branch-law} exactly in the time
      domain and retain an explicit basis of pump harmonics
      (\cref{sec:pump-basis}).

\item \textbf{Undepleted pump.} Assumed throughout
      \eqref{eq:cme-s}--\eqref{eq:gain-cme}.
      \emph{Solver response:} the multitone formulation of
      \cref{chap:saturation} solves pump and signal simultaneously, with no
      undepleted-pump assumption anywhere.

\item \textbf{Slowly-varying envelope.} Fails wherever the design deliberately
      introduces structure on the wavelength scale --- which is exactly what
      dispersion engineering does.
      \emph{Solver response:} solve the discrete node equations
      \eqref{eq:jtwpa-node} with no continuum limit.

\item \textbf{Perfect periodicity.} Real devices have fabrication scatter,
      deliberate tapers, and boundary effects.
      \emph{Solver response:} per-cell component profiles and independent
      per-component scatter (\cref{sec:profiles}).

\item \textbf{Travelling-wave assumption.} A device with imperfect terminations
      or internal impedance steps supports standing waves, and the very notion
      of a position-dependent envelope $A(x)$ growing monotonically becomes
      questionable.
      \emph{Solver response:} full two-port scattering from the circuit, with
      no assumption that the device is matched.
      \begin{caveat}[one studied device is not travelling-wave]
      For the two-chip device studied in this work, the measured ratio of
      backward to forward wave amplitude is $|V^-|/|V^+| = 0.359$. It is a
      standing-wave resonant structure, not a clean travelling-wave amplifier.
      This is \emph{not} attributable to a port mismatch; it is internal.
      Coupled-mode expressions of the form \eqref{eq:gain-cme} should not be
      applied to it without justification.
      \end{caveat}

\item \textbf{Two or three modes.} The coupled-mode equations track signal and
      idler. A driven nonlinear line supports a whole ladder of sidebands at
      $\ws + m\wp$, and for a strongly pumped device the power leaking into
      $|m|>1$ is not small.
      \emph{Solver response:} the Floquet sideband ladder of
      \cref{chap:floquet}, truncated at $|m|\le S$ with $S$ an explicit
      convergence parameter.
\end{enumerate}

The remainder of this thesis constructs the numerical method that removes these
assumptions, and then --- importantly --- examines how much of the result can
actually be trusted.
